% Evans second edition, section 2.5.
% English source: printed p. 81; PDF page 98 (1-based).
% Chinese section starts checked from chapter opening, PDF 31; mean-value pages checked at PDF 36-37 and 53-55.
% Status: editorial study draft with reconstructed proofs and mathematical annotations.
\section{Problems}
\label{sec:2-5}

\noindent\textit{Source: Evans \S2.5, Problems 1--24, printed pp.~81--86
(PDF pp.~98--103).} Statements are concisely restated, and solutions
are editorial. Smoothness is assumed wherever a displayed differentiation
requires it. Problems 3--5 and 17 are especially useful checks on the
mean-value arguments.

\subsection{Problems 1--8: transport and elliptic averages}
\paragraph{1. Transport with a zeroth-order term.}
Solve $u_t+b\cdot Du+cu=0$, $u(\cdot,0)=g$.
\begin{proof}[Solution]
Along $X(s)=x+(s-t)b$, the quantity $e^{cs}u(X(s),s)$ has derivative
zero. Evaluate at $s=0,t$ to obtain $u(x,t)=e^{-ct}g(x-tb)$.
\end{proof}

\paragraph{2. Rotation invariance of the Laplacian.}
\begin{proof}[Solution]
For an orthogonal matrix $O$, the Hessian of $v(x)=u(Ox)$ is
$D^2v=O^T(D^2u)(Ox)O$. Its trace equals
$\operatorname{tr}((D^2u)(Ox)OO^T)=\Delta u(Ox)$.
Thus a harmonic $u$ gives a harmonic $v$.
\end{proof}

\paragraph{3. The center value for Poisson's equation.}
Derive the center formula for $-\Delta u=f$ on $B_r$, $u=g$ on $\partial B_r$.
\begin{proof}[Solution]
Integrate (M) from zero to $r$:
\[
 S_u(0,r)-u(0)
 =-\frac1{n\omega_n}\int_0^r\rho^{1-n}
                         \int_{B_\rho}f(y)\,dy\,d\rho.
\]
Changing the order gives (P), with the nonnegative kernel
$[|y|^{2-n}-r^{2-n}]/[n(n-2)\omega_n]$ when $n\geq3$.
Replace $S_u(0,r)$ by the boundary average of $g$.
The kernel's singularity is locally integrable, justifying the change
for bounded $f$; signed $f$ is treated by absolute integrability.
\end{proof}

\paragraph{4. A direct maximum-principle proof.}
\begin{proof}[Solution]
For $\varepsilon>0$, $v=u+\varepsilon|x|^2$ satisfies
$\Delta v=2n\varepsilon>0$. At an interior maximum its Hessian would
be negative semidefinite and $\Delta v\leq0$, impossible.
Thus $v$ reaches its maximum on the boundary.
Bound the added quadratic uniformly on bounded $U$ and let
$\varepsilon\downarrow0$.
\end{proof}

\paragraph{5. Subharmonic functions.}
Show the mean inequality and boundary maximum principle for $\Delta v\geq0$;
prove that $\phi(u)$ and $|Du|^2$ are subharmonic for harmonic $u$
and smooth convex $\phi$.
\begin{proof}[Solution]
Formula (M) makes spherical means nondecreasing, and (SB) gives
$v(x)\leq A_v(x,r)$.
For the boundary maximum argument use
$v+\varepsilon|x|^2$, as in Problem 4.
The chain rule gives
$\Delta\phi(u)=\phi''(u)|Du|^2+\phi'(u)\Delta u\geq0$.
Since harmonic functions are smooth, differentiating the PDE is valid:
\[
 \Delta|Du|^2=2|D^2u|^2+2Du\cdot D\Delta u=2|D^2u|^2\geq0.
\]
\end{proof}
\Needspace{45mm}
\PDEAnnotation{凸性与平均相容}{调和函数经凸函数变换后不一定仍调和，却成为次调和。Jensen 不等式给直觉，链式法则中的 $\phi''(u)|Du|^2$ 给出可核验的原因。}

\paragraph{6. A uniform Poisson bound.}
\begin{proof}[Solution]
Enclose $U$ in $B_R(0)$, put $F=\|f\|_\infty$, $G=\|g\|_\infty$, and set
$b(x)=G+F(R^2-|x|^2)/(2n)$.
Then $b\geq G$ on $\partial U$ and $-\Delta b=F$.
The function $u-b$ is subharmonic because $-\Delta(u-b)=f-F\leq0$;
its boundary values are nonpositive. Problem 5 yields $u\leq b$.
Repeat for $-u$, giving
$\|u\|_\infty\leq G+R^2F/(2n)$.
\end{proof}

\paragraph{7. Explicit Harnack bounds on a ball.}
\begin{proof}[Solution]
For positive harmonic $u$ on $B_R$ with continuous boundary values,
use (PB), $R-|x|\leq|x-y|\leq R+|x|$, and
$\int_{\partial B_R}u=n\omega_nR^{n-1}u(0)$:
\[
 R^{n-2}\frac{R-|x|}{(R+|x|)^{n-1}}u(0)
 \leq u(x)\leq
 R^{n-2}\frac{R+|x|}{(R-|x|)^{n-1}}u(0).
\]
If $u$ is only given on the open ball, first use radii $\rho<R$
with $|x|<\rho$, then let $\rho\uparrow R$.
\end{proof}

\paragraph{8. Verify the ball Poisson formula.}
\begin{proof}[Solution]
The full verification following (PB) establishes mass one using the
known solution $1$, harmonicity by direct differentiation, and boundary
convergence by a near/far split. Positivity controls the near error,
and the factor $R^2-|x|^2$ kills the far error.
This proves existence; the maximum principle supplies uniqueness.
\end{proof}

\subsection{Problems 9--17: boundaries and the heat equation}
\paragraph{9. Lipschitz boundary data can give an unbounded normal derivative.}
Let bounded continuous half-space data satisfy $g(z)=|z|$ for $|z|\leq1$,
$n\geq2$. Show the Poisson solution has unbounded gradient near zero.
\begin{proof}[Solution]
The boundary value at zero is zero. For $\lambda>0$,
\[
 \frac{u(\lambda e_n)}{\lambda}
 =\frac2{n\omega_n}\int_{\R^{n-1}}
                \frac{g(z)}{(|z|^2+\lambda^2)^{n/2}}\,dz.
\]
The portion $|z|>1$ is uniformly bounded in absolute value, by boundedness
of $g$ and the integrable radial tail.
On $\lambda<|z|<1$, polar coordinates bound the remaining positive
integral below by $c_n\int_\lambda^1 d\rho/\rho$, diverging.
A bounded gradient would bound this difference quotient, a contradiction.
\end{proof}

\paragraph{10. Odd reflection across a flat zero boundary.}
\begin{proof}[Solution]
First assume $u$ is $C^2$ up to the flat boundary of a half-ball,
harmonic inside, and zero on the flat portion.
Tangential first and second derivatives vanish there.
The PDE then makes $u_{nn}=0$ there.
For the odd extension, tangential derivatives change sign below the
plane, normal first derivatives match, mixed derivatives match, and
normal second derivatives match because their boundary values are zero.
Thus the extension is $C^2$ and harmonic.

For only continuous boundary regularity, take a small full ball centered
on the flat plane. Prescribe its continuous odd-reflected spherical data
and solve by (PB). Uniqueness makes this harmonic extension odd,
hence zero on the plane. On the upper half-ball it and $u$ have the
same full boundary data, so the maximum principle makes them equal.
Its lower half is the desired odd reflection. Such balls cover the interface.
\end{proof}
\Needspace{45mm}
\PDEAnnotation{反射要检查到二阶}{奇延拓的函数值匹配只是第一步。要直接证明 $C^2$，还需用零边界值消去切向二阶导数，再用方程推出法向二阶导数为零。}

\paragraph{11. Kelvin transform.}
Show $\overline u(x)=|x|^{2-n}u(x/|x|^2)$ is harmonic when $u$ is harmonic,
on the inverted domain away from zero.
\begin{proof}[Solution]
Put $r=|x|$, $z=x/r^2$ and $e=x/r$.
Then $Dz=r^{-2}(I-2ee^T)$, so $Dz(Dz)^T=r^{-4}I$,
and $\Delta z_i=-2(n-2)x_i/r^4$.
Also $\Delta r^{2-n}=0$ off zero.
In the Laplacian product and chain rule, the first-derivative terms cancel:
the contribution from $r^{2-n}\Delta z_i$ is opposite to that from
$2D(r^{2-n})\cdot Dz_i$. The remaining Hessian term is
\[
 \Delta_x\big[r^{2-n}u(z)\big]=r^{-n-2}(\Delta u)(z).
\]
This proves the claim, including $n=2$ with prefactor one.
\end{proof}

\paragraph{12. Parabolic scaling and its infinitesimal generator.}
\begin{proof}[Solution]
For $\lambda>0$, define $u_\lambda(x,t)=u(\lambda x,\lambda^2t)$. The chain rule gives
\[
 Lu_\lambda=\lambda^2(Lu)(\lambda x,\lambda^2t)=0.
\]
Differentiate this identity in $\lambda$ at one.
The result is $L(x\cdot Du+2tu_t)=0$.
Nonzero negative $\lambda$ also gives a solution by spatial reflection;
$\lambda=0$ requires data at time zero and is not needed for the argument.
\end{proof}

\paragraph{13. A self-similar profile.}
For $u(x,t)=v(x/\sqrt t)$ in one space dimension, derive the kernel.
\begin{proof}[Solution]
For $z=x/\sqrt t$, $u_t=-zv'/(2t)$ and $u_{xx}=v''/t$.
Thus $v''+(z/2)v'=0$, so $v'=c e^{-z^2/4}$ and
$v(z)=c\int_0^z e^{-q^2/4}\,dq+d$.
Choose $c=(4\pi)^{-1/2}$ and $d=1/2$ to obtain initial step data.
Differentiating in $x$ gives
$u_x=(4\pi t)^{-1/2}e^{-x^2/(4t)}=K(x,t)$.
The distributional derivative of the initial step is $\delta_0$,
so this construction indeed produces the fundamental solution.
\end{proof}

\paragraph{14. Heat flow with reaction and forcing.}
\begin{proof}[Solution]
For $u_t-\Delta u+cu=f$, let $v=e^{ct}u$.
Then $Lv=e^{ct}f$ and $v(\cdot,0)=g$. Duhamel gives
\[
 u(x,t)=e^{-ct}(K_t*g)(x)
       +\int_0^te^{-c(t-s)}(K_{t-s}*f(\cdot,s))(x)\,ds.
\]
\end{proof}

\paragraph{15. Heat equation on a half-line with boundary forcing.}
For zero initial data and boundary value $g(t)$, $g(0)=0$, derive
\[
 u(x,t)=\int_0^t
 \frac{x\,e^{-x^2/[4(t-s)]}}{\sqrt{4\pi}(t-s)^{3/2}}g(s)\,ds,\qquad x>0.
\]
\begin{proof}[Solution]
The zero-boundary half-line heat Green kernel is
$K(x-y,\tau)-K(x+y,\tau)$.
Its inward boundary flux gives the Poisson kernel
$P(x,\tau)=-2K_x(x,\tau)
=x e^{-x^2/(4\tau)}/(\sqrt{4\pi}\tau^{3/2})$.
It satisfies the heat equation for $x,\tau>0$.
The substitution $z=x/(2\sqrt\tau)$ shows
$\int_0^\infty P(x,\tau)\,d\tau=1$.
As $x\downarrow0$, the mass on $\tau>\delta$ tends to zero.
Thus $P(x,\cdot)$ is an approximate identity at time zero.
After extending $g$ by zero to negative times, its convolution tends
to $g(t)$; compatibility $g(0)=0$ makes the corner continuous.
For fixed $x>0$, the exponential kills all singular powers at
$\tau=0$, justifying differentiation and the PDE.
The time interval shrinks to zero at the initial limit, proving the
zero initial datum.
\end{proof}

\paragraph{16. Direct parabolic maximum principle.}
\begin{proof}[Solution]
The argument $v=u-\varepsilon t$ in \S2.3.3 supplies the proof.
At a top-time maximum the left derivative is nonnegative; one does not
need to pretend that the terminal point is a two-sided time maximum.
\end{proof}

\paragraph{17. Subsolutions and convex transformations.}
\begin{proof}[Solution]
For $Lv\leq0$, (HR) and its limit at zero imply $v(x,t)\leq M_r[v]$.
The perturbed maximum argument from Problem 16 applies to subsolutions,
giving the boundary maximum principle on bounded cylinders.
If $\phi$ is smooth and convex and $Lu=0$, then
\[
 L\phi(u)=\phi'(u)Lu-\phi''(u)|Du|^2\leq0.
\]
For a smooth caloric $u$, spatial and time derivatives are caloric,
and $L(w^2)=-2|Dw|^2$ whenever $Lw=0$.
Sum for $w=\partial_i u$ and $w=u_t$:
\[
 L(|Du|^2+u_t^2)=-2|D^2u|^2-2|Du_t|^2\leq0.
\]
\end{proof}
\Needspace{45mm}
\PDEAnnotation{椭圆与热的凸性公式}{椭圆侧是 $\Delta\phi(u)=\phi''(u)|Du|^2\geq0$；热算子侧是 $L\phi(u)=-\phi''(u)|Du|^2\leq0$。符号不同，得到的平均不等式方向一致。}

\subsection{Problems 18--24: waves, systems, and energy}
\paragraph{18. Stokes' rule.}
If $u$ has wave data $(0,h)$, show $v=u_t$ has data $(h,0)$.
\begin{proof}[Solution]
Time differentiation commutes with $\partial_t^2-\Delta$, so $v$ solves
the same PDE. At zero, $v=h$ and
$v_t=u_{tt}=\Delta u=\Delta0=0$. This is also the reason the displacement
part of the solution is a time derivative of the velocity solution operator.
\end{proof}

\paragraph{19. Null coordinates and one-way waves.}
\begin{proof}[Solution]
The equation $u_{xy}=0$ first makes $u_x$ independent of $y$;
integrating gives $u=F(x)+G(y)$.
For $\xi=x+t$, $\eta=x-t$, the wave operator is
$-4\partial_\xi\partial_\eta$, so
$u=F(x+t)+G(x-t)$. The initial conditions give
$F'=(g'+h)/2$, $G'=(g'-h)/2$, producing (DA).
A purely right-moving profile $G(x-t)$ requires $h=-g'$.
A purely left-moving profile $F(x+t)$ requires $h=g'$.
Additive constants can be allocated to either profile without changing
these derivative conditions.
\end{proof}

\paragraph{20. Undistorted radial profiles.}
For nonzero attenuation $\alpha(r)$ and nonnegative delay $\beta(r)$,
$\beta(0)=0$, suppose $\alpha(r)\phi(t-\beta(r))$ solves the wave
equation for every smooth profile $\phi$.
\begin{proof}[Solution]
Insert the form into $u_{tt}-u_{rr}-(n-1)u_r/r$.
Independence of $\phi,\phi',\phi''$ gives
\[
 \alpha(1-\beta'^2)=0,\qquad
 2\alpha'\beta'+\alpha\beta''+\frac{n-1}r\alpha\beta'=0,
\]
\[
 \alpha''+\frac{n-1}r\alpha'=0.
\]
For a nontrivial smooth attenuation, the first and the nonnegative-delay
condition give $\beta(r)=r$; the middle equation gives
$\alpha=Cr^{-(n-1)/2}$.
Substitution into the last equation yields
$(n-1)(n-3)C=0$. Thus $n=1$, $\alpha=C$, or $n=3$, $\alpha=C/r$.
The identically zero attenuation is a trivial exception and has to be
excluded to make the dimension claim true.
\end{proof}
\Needspace{45mm}
\PDEAnnotation{量词“任意波形”很强}{把 $\phi,\phi',\phi''$ 的系数分别置零，依赖于要求对所有光滑 $\phi$ 成立。若只要求某一个特定波形，不能推出相同的维数限制。}

\paragraph{21. Maxwell and elastic waves.}
\begin{proof}[Solution]
Differentiate $E_t=\operatorname{curl}B$ and use
$B_t=-\operatorname{curl}E$:
\[
 E_{tt}=-\operatorname{curl}\operatorname{curl}E
       =-\nabla\operatorname{div}E+\Delta E=\Delta E.
\]
Likewise $B_{tt}=\Delta B$, using the divergence constraints.
For elasticity,
$u_{tt}-\mu\Delta u-(\lambda+\mu)\nabla\operatorname{div}u=0$.
Taking divergence and curl respectively gives
\[
 w_{tt}-(\lambda+2\mu)\Delta w=0,\qquad
 \mathbf w_{tt}-\mu\Delta\mathbf w=0,
\]
where $w=\operatorname{div}u$ and $\mathbf w=\operatorname{curl}u$.
When $\mu>0$ and $\lambda+2\mu>0$, the longitudinal and transverse
speeds are $\sqrt{\lambda+2\mu}$ and $\sqrt\mu$.
\end{proof}

\paragraph{22. Telegraph equation from velocity switching.}
Given $u_t+u_x=d(v-u)$ and $v_t-v_x=d(u-v)$, show each component satisfies
$w_{tt}+2dw_t-w_{xx}=0$.
\begin{proof}[Solution]
Let $a=u+v$, $b=u-v$. Adding and subtracting gives
$a_t+b_x=0$ and $b_t+a_x=-2db$.
Differentiate the first in time:
$a_{tt}=-b_{xt}=a_{xx}+2db_x=a_{xx}-2da_t$.
Differentiate the second in time and use $a_{xt}=-b_{xx}$:
$b_{tt}=b_{xx}-2db_t$.
Linear combinations give the same telegraph equation for $u,v$.
\end{proof}

\paragraph{23. A compact space-time source and its outgoing shape.}
In the diamond $S=\{|x|+|t-2|<1\}$, take $f=1$ below $t=x+2$
and $f=-1$ above it, zero elsewhere. Solve with zero wave data
and describe $u$ for $t>3$.
\begin{proof}[Solution]
The one-dimensional Duhamel formula is
\[
 u(x,t)=\frac12\int_0^t
                 \int_{x-(t-s)}^{x+(t-s)}f(y,s)\,dy\,ds.
\]
Set $\alpha=s+y$, $\beta=s-y$.
Then $S$ becomes $1<\alpha<3$, $1<\beta<3$,
and $dy\,ds=\tfrac12d\alpha\,d\beta$.
The sign depends only on $\beta$: it is $+1$ for $\beta<2$,
$-1$ for $\beta>2$.
The backward cone imposes $\alpha\leq t+x$, $\beta\leq t-x$.
Define
\[
 F(a)=\min\{2,\max\{0,a-1\}\},
\]
\[
 H(b)=
 \begin{cases}
 0,&b\leq1,\\
 b-1,&1<b\leq2,\\
 3-b,&2<b<3,\\
 0,&b\geq3.
 \end{cases}
\]
The integrals factor, giving $u(x,t)=\tfrac14F(t+x)H(t-x)$.
If $t>3$ and $H(t-x)\ne0$, then $t+x=2t-(t-x)>3$,
so $F(t+x)=2$. Therefore
\[
 u(x,t)=\tfrac12H(t-x)\quad(t>3).
\]
It is a right-moving triangular pulse of height $1/2$, supported on
$t-3<x<t-1$, with peak at $x=t-2$. The sign-balanced source cancels
the other outgoing pulse.
This piecewise forcing is interpreted through the integral formula;
across its jump lines the PDE is read distributionally.
\end{proof}
\Needspace{45mm}
\PDEAnnotation{特征坐标把菱形变方形}{$\alpha=s+y,\beta=s-y$ 同时把波锥不等式变成坐标上界。二重积分才能分解成两个一维截断积分；漏掉 Jacobian 的 $1/2$ 会让最终脉冲高度错一倍。}

\paragraph{24. Equipartition in one dimension.}
For compactly supported smooth $g,h$, show total wave energy is constant
and kinetic and potential energies eventually coincide.
\begin{proof}[Solution]
Write $u=F(x+t)+G(x-t)$ as in Problem 19.
Although the primitives $F,G$ need not be compactly supported,
their derivatives $F'=(g'+h)/2$, $G'=(g'-h)/2$ are.
Since $u_t=F'-G'$ and $u_x=F'+G'$,
\[
 k(t)+p(t)=\int_{\R}\bigl(F'(x+t)^2+G'(x-t)^2\bigr)\,dx
\]
is independent of $t$ by translation invariance. Also
\[
 k(t)-p(t)=-2\int_{\R} F'(x+t)G'(x-t)\,dx.
\]
For sufficiently large $t$ the translated compact supports are disjoint,
so this integral is zero. If the initial supports lie in $[-R,R]$,
every $t>R$ suffices.
\end{proof}
